\documentclass[a4paper,12pt]{article}

\usepackage[T2A]{fontenc}
\usepackage[utf8]{inputenc}
\usepackage[russian]{babel}

\usepackage{amsmath,amssymb,amsthm,mathtools}
\usepackage{helvet}
\renewcommand{\familydefault}{\sfdefault}
\usepackage{icomma}
\usepackage{geometry}
\usepackage{microtype}
\usepackage{tikz}
\usepackage{tkz-euclide}
\usepackage{pgfplots}
\pgfplotsset{compat=1.18}
\usepackage{tabularx,booktabs,longtable}
\usepackage{enumitem}
\usepackage{hyperref}
\usepackage{parskip}
\usepackage{fancyhdr}
\usepackage{setspace}
\usepackage{xcolor}
\usepackage{array}

\geometry{margin=2.2cm}
\onehalfspacing
\hypersetup{hidelinks}

\definecolor{accent}{HTML}{008080}
\definecolor{darkaccent}{HTML}{004D4D}
\definecolor{textgray}{HTML}{333333}

\pagestyle{fancy}
\fancyhf{}
\fancyhead[L]{\textcolor{darkaccent}{Чек-лист по стереометрии}}
\fancyhead[R]{\textcolor{darkaccent}{Призмы}}
\fancyfoot[C]{\thepage}

\tkzSetUpLine[line width=0.9pt, color=accent]
\tkzSetUpPoint[color=accent, fill=accent, size=3]
\tkzSetUpLabel[color=black, font=\small]

\setlist[itemize]{leftmargin=1.8em}
\setlist[enumerate]{leftmargin=2em}

\newcommand{\ds}{\displaystyle}
\newcommand{\Area}{S}
\newcommand{\Vol}{V}
\newcommand{\per}{P}

\begin{document}

\begin{titlepage}
\thispagestyle{empty}
\begin{center}
\vspace*{2.8cm}

{\Large\bfseries\textcolor{darkaccent}{Чек-лист}}\\[0.6cm]
{\LARGE\bfseries Призмы: объём, площадь поверхности,}\\[0.2cm]
{\LARGE\bfseries работа с основаниями и отношения величин}

\vspace{2.2cm}

{\large Лёвин Артём Александрович эксклюзивно для Яны Плиско}

\vspace{0.9cm}

{\large \today}

\vfill

\begin{tikzpicture}[x=1cm,y=1cm]
\draw[accent, line width=1.2pt]
(0,0) .. controls (3.5,1.1) and (7,1.1) .. (10.5,0.25)
      .. controls (12.2,-0.1) and (14.2,-0.1) .. (16,0.5);
\draw[darkaccent, line width=0.5pt]
(0,0.35) .. controls (3.7,1.35) and (7.3,1.35) .. (10.7,0.55)
           .. controls (12.1,0.22) and (14,0.22) .. (16,0.8);
\end{tikzpicture}

\end{center}
\end{titlepage}

\section*{Основной чек-лист}

\subsection*{1. Что обязательно помнить о призме}

\begin{itemize}
    \item \textbf{Призма} --- многогранник, у которого две грани являются \textbf{равными и параллельными многоугольниками}. Эти грани называются \textbf{основаниями}.
    \item Все остальные грани --- \textbf{боковые}. У прямой призмы они являются прямоугольниками.
    \item \textbf{Высота призмы} --- расстояние между плоскостями оснований.
    \item \textbf{Боковое ребро} совпадает с высотой \textbf{только у прямой призмы}.
    \item \textbf{Правильная призма} в школьных задачах --- это \textbf{прямая призма}, в основании которой лежит \textbf{правильный многоугольник}.
\end{itemize}

\subsection*{2. Главные формулы}

\begin{itemize}
    \item \textbf{Объём призмы:}
    \[
    \Vol=\Area_{\text{осн}}\cdot h.
    \]
    
    \item \textbf{Площадь полной поверхности:}
    \[
    \Area_{\text{полн}}=2\Area_{\text{осн}}+\Area_{\text{бок}}.
    \]
    
    \item \textbf{Для прямой призмы площадь боковой поверхности:}
    \[
    \Area_{\text{бок}}=\per_{\text{осн}}\cdot h.
    \]
    
    \item Здесь:
    \[
    \Area_{\text{осн}}=\text{площадь основания},\qquad
    \per_{\text{осн}}=\text{периметр основания},\qquad
    h=\text{высота}.
    \]
\end{itemize}

\subsection*{3. Площади оснований, которые встречаются чаще всего}

\begin{itemize}
    \item \textbf{Прямоугольный треугольник с катетами $a$ и $b$:}
    \[
    \Area=\frac{ab}{2}.
    \]
    
    \item \textbf{Равносторонний треугольник со стороной $a$:}
    \[
    \Area=\frac{a^{2}\sqrt{3}}{4}.
    \]
    
    \item \textbf{Ромб с диагоналями $d_1$ и $d_2$:}
    \[
    \Area=\frac{d_1d_2}{2}.
    \]
    
    \item \textbf{Правильный шестиугольник со стороной $a$:}
    \[
    \Area=\frac{3\sqrt{3}}{2}a^2.
    \]
\end{itemize}

\subsection*{4. Полезные факты про основания}

\begin{itemize}
    \item Если основание --- \textbf{ромб}, то его диагонали:
    \begin{itemize}
        \item пересекаются и делятся пополам;
        \item взаимно перпендикулярны.
    \end{itemize}
    
    \item Если основание --- \textbf{равносторонний треугольник}, то для площади удобно сразу применять формулу
    \[
    \Area=\frac{a^{2}\sqrt{3}}{4}.
    \]
    
    \item Если в основании есть \textbf{прямоугольный треугольник}, не забывайте про теорему Пифагора:
    \[
    c^2=a^2+b^2.
    \]
    
    \item \textbf{Средняя линия треугольника} параллельна третьей стороне и равна половине этой стороны.
\end{itemize}

\subsection*{5. Базовый алгоритм для задач на объём}

\begin{enumerate}
    \item Определите, какая фигура лежит в основании.
    \item Найдите площадь основания по подходящей формуле.
    \item Найдите высоту призмы.
    \item Подставьте в формулу
    \[
    \Vol=\Area_{\text{осн}}\cdot h.
    \]
\end{enumerate}

\textbf{Мини-пример.} У прямой призмы в основании прямоугольный треугольник с катетами $3$ и $5$, а высота равна $4$. Тогда
\[
\Area_{\text{осн}}=\frac{3\cdot 5}{2}=\frac{15}{2},
\qquad
\Vol=\frac{15}{2}\cdot 4=30.
\]

\subsection*{6. Базовый алгоритм для задач на площадь полной поверхности}

\begin{enumerate}
    \item Найдите площадь одного основания.
    \item Найдите периметр основания.
    \item Найдите площадь боковой поверхности по формуле
    \[
    \Area_{\text{бок}}=\per_{\text{осн}}\cdot h
    \]
    (только для прямой призмы).
    \item Используйте формулу
    \[
    \Area_{\text{полн}}=2\Area_{\text{осн}}+\Area_{\text{бок}}.
    \]
\end{enumerate}

\textbf{Мини-пример.} У прямой призмы в основании ромб с диагоналями $3$ и $4$, а высота равна $5$.
\[
\Area_{\text{осн}}=\frac{3\cdot 4}{2}=6.
\]
Половины диагоналей равны $1{,}5$ и $2$, значит сторона ромба:
\[
a=\sqrt{1{,}5^2+2^2}=\sqrt{2{,}25+4}=2{,}5.
\]
Тогда
\[
\per_{\text{осн}}=4a=10,
\qquad
\Area_{\text{бок}}=10\cdot 5=50,
\]
\[
\Area_{\text{полн}}=2\cdot 6+50=62.
\]

\subsection*{7. Метод отношений: когда удобно делить одну формулу на другую}

Этот приём особенно полезен, когда:
\begin{itemize}
    \item просят сравнить \textbf{объёмы};
    \item просят сравнить \textbf{площади};
    \item одна фигура получена из другой изменением стороны основания и/или высоты;
    \item в задаче есть \textbf{средние линии}, \textbf{увеличение в несколько раз}, \textbf{уменьшение в несколько раз}.
\end{itemize}

\textbf{Идея:} записать формулу для исходной фигуры и для новой фигуры, затем \textbf{разделить одно на другое}.

\paragraph{Для объёма}
\[
\frac{\Vol_2}{\Vol_1}
=
\frac{\Area_{\text{осн},2}\cdot h_2}{\Area_{\text{осн},1}\cdot h_1}.
\]

\paragraph{Для площади боковой поверхности прямой призмы}
\[
\frac{\Area_{\text{бок},2}}{\Area_{\text{бок},1}}
=
\frac{\per_{\text{осн},2}\cdot h_2}{\per_{\text{осн},1}\cdot h_1}.
\]

\subsection*{8. Очень важные частные случаи масштабирования}

\begin{itemize}
    \item Если \textbf{все линейные размеры основания} увеличили в $k$ раз, то
    \[
    \Area_{\text{осн}}
    \text{ увеличится в } k^2 \text{ раз.}
    \]
    
    \item Если \textbf{периметр} подобной фигуры увеличили в $k$ раз, то он увеличится просто в $k$ раз.
    
    \item Если у призмы:
    \begin{itemize}
        \item сторона основания изменилась в $k$ раз;
        \item высота изменилась в $m$ раз,
    \end{itemize}
    то
    \[
    \Vol \text{ изменится в } k^2m \text{ раз.}
    \]
    
    \item Для площади боковой поверхности прямой призмы:
    \[
    \Area_{\text{бок}} \text{ изменится в } km \text{ раз,}
    \]
    потому что
    \[
    \Area_{\text{бок}}=\per_{\text{осн}}\cdot h,
    \]
    а периметр --- величина линейная.
\end{itemize}

\textbf{Мини-пример.} Если сторону равностороннего основания увеличили в $3$ раза, а высоту уменьшили в $2$ раза, то объём изменится в
\[
3^2\cdot \frac12=\frac92
\]
раза.

\subsection*{9. Схема к задаче про ромб в основании}

Ниже показан полезный чертёж: диагонали ромба перпендикулярны и делят друг друга пополам, поэтому сторону ромба удобно искать по теореме Пифагора.

\begin{center}
\begin{tikzpicture}[scale=1.1]
\tkzDefPoint(0,0){O}
\tkzDefPoint(-3,0){A}
\tkzDefPoint(0,2){B}
\tkzDefPoint(3,0){C}
\tkzDefPoint(0,-2){D}

\tkzDrawPolygon[thick,color=accent](A,B,C,D)
\tkzDrawSegments[dashed,color=darkaccent](A,C B,D)
\tkzMarkRightAngle[size=0.25](A,O,B)

\tkzDrawPoints(O,A,B,C,D)
\tkzLabelPoints[below](A,C,D)
\tkzLabelPoints[above](B)
\tkzLabelPoints[above right](O)

\node[below left] at (-1.5,0) {\small $\dfrac{d_1}{2}$};
\node[above right] at (0,1) {\small $\dfrac{d_2}{2}$};
\end{tikzpicture}
\end{center}

Следовательно,
\[
a=\sqrt{\left(\frac{d_1}{2}\right)^2+\left(\frac{d_2}{2}\right)^2}.
\]

\subsection*{10. Типичные ошибки}

\begin{itemize}
    \item Путать \textbf{правильную} и \textbf{прямую} призму.
    \item Автоматически считать боковое ребро высотой, когда призма не названа прямой.
    \item В формуле объёма подставлять не площадь основания, а что-то линейное.
    \item В формуле площади боковой поверхности забывать, что нужен \textbf{периметр основания}, а не площадь.
    \item Для ромба ошибочно брать сторону равной половине диагонали.
    \item Для равностороннего треугольника забывать специальную формулу
    \[
    \Area=\frac{a^2\sqrt{3}}{4}.
    \]
    \item При задачах на отношения забывать, что площадь меняется \textbf{квадратично}, а периметр --- \textbf{линейно}.
    \item Не замечать среднюю линию треугольника и, как следствие, не использовать отношение $1:2$ для соответствующих сторон.
\end{itemize}

\subsection*{11. Что делать, если задача не решается сразу}

\begin{enumerate}
    \item Выпишите, \textbf{что именно является основанием}.
    \item Найдите или выразите \textbf{площадь основания}.
    \item Найдите \textbf{высоту}.
    \item Если просят поверхность, найдите \textbf{периметр основания}.
    \item Если в основании есть ромб, треугольник, средняя линия, диагонали --- почти наверняка поможет \textbf{планиметрия}.
    \item Если нужно сравнить две призмы --- попробуйте \textbf{метод деления формул}.
\end{enumerate}

\subsection*{12. Короткий итог по теме}

\begin{itemize}
    \item Для объёма главное:
    \[
    \Vol=\Area_{\text{осн}}\cdot h.
    \]
    \item Для полной поверхности прямой призмы главное:
    \[
    \Area_{\text{полн}}=2\Area_{\text{осн}}+\per_{\text{осн}}\cdot h.
    \]
    \item Самое важное в большинстве задач --- не стереометрия сама по себе, а \textbf{умение работать с фигурой в основании}.
\end{itemize}

\newpage

\section*{Тренировочный блок}

\textbf{Простые задачи}

\begin{enumerate}[label=\arabic*.]
    \item Найдите объём правильной треугольной призмы, если все её рёбра равны $\sqrt{3}$.
    
    \item Основанием прямой треугольной призмы служит прямоугольный треугольник с катетами $3$ и $5$. Объём призмы равен $30$. Найдите боковое ребро призмы.
    
    \item Найдите объём прямой призмы, если в основании лежит равносторонний треугольник со стороной $4$, а высота призмы равна $6$.
\end{enumerate}

\textbf{Средние задачи}

\begin{enumerate}[label=\arabic*.,resume]
    \item Найдите площадь полной поверхности прямой призмы, в основании которой лежит ромб с диагоналями $3$ и $4$, если боковое ребро равно $5$.
    
    \item В основании прямой призмы лежит прямоугольный треугольник с катетами $6$ и $8$. Высота призмы равна $10$. Найдите площадь полной поверхности призмы.
    
    \item В основании прямой призмы лежит правильный шестиугольник со стороной $2$. Высота призмы равна $7$. Найдите объём призмы.
\end{enumerate}

\textbf{Задачи выше среднего}

\begin{enumerate}[label=\arabic*.,resume]
    \item Через среднюю линию основания треугольной призмы проведена плоскость, параллельная боковым рёбрам. Площадь боковой поверхности исходной призмы равна $12$. Найдите площадь боковой поверхности отсечённой призмы.
    
    \item У правильной треугольной призмы объём равен $6$. Сторону основания новой призмы увеличили в $3$ раза, а высоту уменьшили в $2$ раза. Найдите объём новой призмы.
    
    \item В основании прямой призмы лежит ромб с диагоналями $6$ и $8$. Площадь боковой поверхности призмы равна $140$. Найдите высоту призмы.
\end{enumerate}

\textbf{Сложная задача}

\begin{enumerate}[label=\arabic*.,resume]
    \item В основании прямой треугольной призмы лежит треугольник со сторонами $13$, $14$ и $15$. Площадь полной поверхности призмы равна $454$. Найдите объём призмы.
\end{enumerate}

\newpage

\section*{Ответы}

\renewcommand{\arraystretch}{1.35}

\begin{table}[h!]
\centering
\caption{Ответы к тренировочному блоку}
\begin{tabularx}{\linewidth}{>{\bfseries}c X}
\toprule
№ & Ответ \\
\midrule
1 & $\ds \Vol=\frac{9}{4}$. \\

2 & $\ds h=4$. \\

3 & $\ds \Vol=\Area_{\text{осн}}\cdot h=\frac{4^2\sqrt{3}}{4}\cdot 6=24\sqrt{3}$. \\

4 & $\ds \Area_{\text{полн}}=62$. \\

5 & Основание: $\ds \Area_{\text{осн}}=\frac{6\cdot 8}{2}=24$, гипотенуза $10$, периметр основания $24$, боковая поверхность $240$, полная поверхность $\ds 288$. \\

6 & $\ds \Vol=\frac{3\sqrt{3}}{2}\cdot 2^2\cdot 7=42\sqrt{3}$. \\

7 & $\ds 6$. \\

8 & $\ds 27$. \\

9 & Половины диагоналей равны $3$ и $4$, сторона ромба $5$, периметр основания $20$, поэтому $\ds h=\frac{140}{20}=7$. \\

10 & По формуле Герона для треугольника со сторонами $13$, $14$, $15$:
$\ds p=\frac{13+14+15}{2}=21$,
$\ds \Area_{\text{осн}}=\sqrt{21\cdot 8\cdot 7\cdot 6}=84$.
Периметр основания равен $42$, значит
$\ds 2\cdot 84+42h=454$,
откуда $\ds h=\frac{286}{42}=\frac{143}{21}$.
Тогда
$\ds \Vol=84\cdot \frac{143}{21}=572$. \\
\bottomrule
\end{tabularx}
\end{table}

\end{document}